\documentclass[11pt]{article}

\usepackage[margin=1in]{geometry}
\usepackage{amsmath,amssymb,amsthm}
\usepackage{graphicx}
\usepackage{booktabs}
\usepackage{hyperref}
\usepackage{siunitx}
\usepackage{cases}

\newtheorem{definition}{Definition}
\newtheorem{proposition}{Proposition}
\newtheorem{remark}{Remark}

\title{A Mathematical Model for Variable-Geometry Helical\\
Compression Springs with Contact-Aware Stiffening}
\author{Derived from the \texttt{CompressionSpringGeneral} implementation}
\date{}

\begin{document}
\maketitle

\begin{abstract}
We formalize the mathematical model underlying a computational class for
general (possibly variable-diameter) helical compression springs. The model
covers: the local spring index and Wahl stress-correction factor for a
spring whose mean coil diameter $D(z)$ varies along its axis; a geometric
criterion distinguishing \emph{telescoping} (nesting) springs from
conventional stacking springs, used to predict solid length; a
piecewise/incremental contact model of progressive compression that
captures the stiffening produced by coil-to-coil and end-plate contact;
the resulting position-dependent load, shear stress, and outer/inner
diameter fields; and a Goodman fatigue-diagram formulation for
alternating torsional loading. The presentation follows the structure of
the reference implementation, translating each computational routine into
its governing equations.
\end{abstract}

\tableofcontents

\section{Notation and basic geometry}

\begin{table}[h!]
\centering
\begin{tabular}{cl}
\toprule
Symbol & Meaning \\
\midrule
$d$ & wire diameter \\
$D(z)$ & local mean coil diameter at axial coordinate $z$ \\
$L_0$ & free (unloaded) length of the spring \\
$L_s$ & solid length \\
$n$ & number of active coils \\
$k$ & spring constant (stiffness) \\
$C(z)$ & local spring index, $C(z) = D(z)/d$ \\
$K_W(z)$ & Wahl factor at position $z$ \\
$\delta$ & travel (deflection) from free length \\
$F$ & axial load \\
$\tau$ & maximum shear stress in the wire \\
$\nu$ & Poisson's ratio of the wire material \\
$N_c$ & number of fatigue cycles \\
$\sigma_{\max},\sigma_{\min}$ & maximum/minimum alternating shear stress \\
\bottomrule
\end{tabular}
\caption{Principal symbols.}
\end{table}

A spring occupying the free-length interval $z\in[0,L_0]$ is described by
a (possibly non-constant) mean-diameter function $D(z)$, generalizing the
constant-diameter cylindrical spring to conical, barrel, or other
variable-geometry profiles. Positions are addressed either by their
absolute coordinate $z$ (``position'', measured from the free/loaded end)
or by their travel
\begin{equation}
\delta = L_0 - z ,
\label{eq:travel}
\end{equation}
i.e.\ the amount by which the spring has been compressed from its free
length.

\section{Local spring index and the Wahl factor}

\begin{definition}[Local spring index]
At axial position $z$,
\begin{equation}
C(z) = \frac{D(z)}{d}.
\end{equation}
\end{definition}

The Wahl factor corrects the elementary torsion formula for a helical
spring to account for direct shear and coil curvature effects. For local
spring index $C(z)$,
\begin{equation}
K_W(z) = \frac{4C(z)-1}{4C(z)-4} + \frac{0.615}{C(z)} .
\label{eq:wahl}
\end{equation}

\begin{remark}
Equation~\eqref{eq:wahl} is evaluated pointwise using the \emph{local}
index $C(z)$ rather than a single global index, so that a
variable-diameter spring receives a position-dependent correction factor
$K_W(z)$ instead of one constant value.
\end{remark}

Springs are further classified into qualitative severity bands according
to threshold values $c_{\mathrm{r}} $ and $c_{\mathrm{o}}$ on the spring
index:
\begin{equation}
\text{category}(z) =
\begin{cases}
\text{``red''}, & C(z) < c_{\mathrm r} \\[2pt]
\text{``orange''}, & c_{\mathrm r} \le C(z) < c_{\mathrm o} \\[2pt]
\text{``green''}, & C(z) \ge c_{\mathrm o}
\end{cases}
\end{equation}
where $c_{\mathrm r} < c_{\mathrm o}$ are material/design constants,
flagging low spring indices (tightly wound coils) as more severe because
the Wahl correction grows sharply as $C\to 1^{+}$.

\section{Solid length: stacking versus telescoping}

When a compression spring is closed completely, its solid length $L_s$
depends on whether successive coils can nest \emph{inside} one another
(telescoping, possible for springs whose diameter changes rapidly enough
along their length, e.g.\ strongly conical springs) or must stack
\emph{on top of} one another (the conventional case for near-cylindrical
springs).

\begin{definition}[Telescoping criterion]
Let
\begin{equation}
D_0 = D(0), \qquad D_1 = D(L_0)
\end{equation}
be the mean diameters at the two ends of the free length. The spring is
classified as telescoping when the total diametral change exceeds the
aggregate wire thickness that would otherwise be needed to stack every
coil:
\begin{equation}
\lvert D_0 - D_1 \rvert \;\ge\; n\,d .
\label{eq:telescope-criterion}
\end{equation}
\end{definition}

\begin{proposition}[Solid length]
\begin{equation}
L_s =
\begin{cases}
d, & \text{if condition \eqref{eq:telescope-criterion} holds
(ideal telescoping limit),}\\[4pt]
n\,d, & \text{otherwise (conventional stacking).}
\end{cases}
\label{eq:solid-length}
\end{equation}
\end{proposition}

The first branch is the limiting case in which every coil can nest fully
inside its neighbor, so the compressed stack collapses to a single wire
diameter; the second is the ordinary case of $n$ coils of diameter $d$
stacking without interference.

\section{Contact-aware progressive compression}

Rather than assuming a single constant stiffness $k$ over the full travel
range $\delta \in [0, L_0 - L_s]$, the load--deflection relation is
obtained from an incremental simulation that can capture stiffening
transitions caused by (i) coil-to-coil contact and (ii) the moving plate
or floor approaching the solid stack.

\begin{definition}[Progressive compression map]
Let $\delta_{\max} = L_0$ and partition $[0,\delta_{\max}]$ into $M$
increments $0=\delta_0 <\delta_1<\dots<\delta_M=\delta_{\max}$. The
simulation produces a discrete force history
\begin{equation}
\big\{(\delta_i, F_i)\big\}_{i=0}^{M},
\end{equation}
in which each $F_i$ is obtained from the instantaneous effective
stiffness $k_{\mathrm{eff}}(\delta_i)$, itself a function of which coils
are currently in contact:
\begin{equation}
F_i = F_{i-1} + k_{\mathrm{eff}}(\delta_{i-1})\,(\delta_i - \delta_{i-1}),
\qquad F_0 = 0 .
\label{eq:incremental-force}
\end{equation}
\end{definition}

The continuous load--travel curve used elsewhere in the model,
\begin{equation}
F(\delta) \approx \operatorname{interp}\big(\delta;\;\{\delta_i\},\{F_i\}\big),
\label{eq:interp}
\end{equation}
is then recovered from the discrete history by piecewise-linear
interpolation, and it is this curve — rather than the naive linear
relation $F = k\delta$ — that supplies the load at any requested
position.

\begin{remark}
Because the simulation is expensive relative to interpolation, it is
evaluated once over the full travel range and reused: both the
load-at-position query (Section~\ref{sec:load-position}) and the
diagnostic load--travel/load--position plots draw on the same discrete
history $\{(\delta_i,F_i)\}$.
\end{remark}

\section{Load, stress, and diameter at a prescribed position}
\label{sec:load-position}

\subsection{Load at a position}

For a requested position $z$ with travel $\delta = L_0 - z$
(Eq.~\eqref{eq:travel}), the load is read from the interpolated
progressive-compression curve and signed by the loading class
$s\in\{-1,0,1\}$ (with $s=1$ for compression):
\begin{equation}
F(z) = s\cdot F\big(L_0 - z\big).
\label{eq:load-at-position}
\end{equation}

\subsection{Shear stress at a position}

The corrected torsional shear stress in the wire uses the elementary
torsion-of-a-curved-bar formula for a round wire of diameter $d$, scaled
by the local Wahl factor $K_W(z)$ of Eq.~\eqref{eq:wahl} and the local
mean diameter $D(z)$:
\begin{equation}
\tau(z) = \frac{8\,D(z)\,F(z)}{\pi d^{3}}\,K_W(z).
\label{eq:stress}
\end{equation}

\subsection{Outer and inner diameter at a position}

The outer diameter accounts for both the wire cross-section and a
Poisson-effect radial growth of the coil as the wire is compressed
axially between the free position and the current position:
\begin{equation}
D_{\mathrm{out}}(z) = D(z) + d + D(z)\,\nu\,\frac{L_0 - z}{L_0},
\label{eq:outer-diameter}
\end{equation}
\begin{equation}
D_{\mathrm{in}}(z) = \max\!\big(D_{\mathrm{out}}(z) - 2d,\; 0\big).
\label{eq:inner-diameter}
\end{equation}
The linear factor $(L_0-z)/L_0 \in [0,1]$ scales the Poisson bulge from
zero at the free end of travel ($z=L_0$) up to its full value as the
spring approaches its most compressed recorded position ($z\to 0$).

\subsection{Admissibility of a position}

A requested position $z$ (equivalently a load-position record) is
physically admissible only if it does not correspond to over-closing the
spring beyond its solid length:
\begin{equation}
z \;\ge\; L_s ,
\end{equation}
with $L_s$ given by Eq.~\eqref{eq:solid-length}; violations are rejected.

\section{Extremal quantities over recorded positions}

Given a finite set of recorded load positions
$\mathcal{P} = \{z_1,\dots,z_p\}$ with associated stresses
$\tau(z_j)$ and loads $F(z_j)$ from Eqs.~\eqref{eq:stress}
and~\eqref{eq:load-at-position}, the reported extremes are simply
\begin{equation}
\tau_{\max} = \max_{j} \tau(z_j), \qquad
\tau_{\min} = \min_{j} \tau(z_j),
\end{equation}
\begin{equation}
F_{\max} = \max_{j} F(z_j), \qquad
F_{\min} = \min_{j} F(z_j).
\end{equation}
These extrema define the alternating stress state used for fatigue
analysis below.

\section{Fatigue analysis: the Goodman criterion}

For a wire loaded in torsion for $N_c$ cycles, define the mean and
alternating shear stress components from the extremal stresses of the
duty cycle,
\begin{equation}
\tau_m = \frac{\tau_{\max}+\tau_{\min}}{2}, \qquad
\tau_a = \frac{\tau_{\max}-\tau_{\min}}{2}.
\label{eq:mean-alt}
\end{equation}

Let $\tau_e$ denote the torsional endurance limit for $N_c$ cycles
(material- and surface-treatment–dependent, e.g.\ adjusted for shot
peening) and $\tau_y$ the torsional yield/ultimate shear strength. The
modified-Goodman fatigue criterion states that a stress state is safe
when
\begin{equation}
\frac{\tau_a}{\tau_e} + \frac{\tau_m}{\tau_y} \;\le\; 1 .
\label{eq:goodman}
\end{equation}
Equation~\eqref{eq:goodman} defines the Goodman line in the
$(\tau_m,\tau_a)$ plane; the operating point $(\tau_m,\tau_a)$ computed
from Eq.~\eqref{eq:mean-alt} is plotted against this line, and its safety
factor is
\begin{equation}
n_f = \left(\frac{\tau_a}{\tau_e} + \frac{\tau_m}{\tau_y}\right)^{-1}.
\label{eq:safety-factor}
\end{equation}

\section{Geometric development of the helical centerline}

For visualization and for consistency with the wire-length calculation,
the coil centerline is parametrized by an angular sweep
$\theta \in [0, 2\pi n]$ mapped to axial coordinate $z(\theta)$, so that
at $N$ sample points $\theta_i$ the three-dimensional centerline is
\begin{equation}
\begin{aligned}
x(\theta_i) &= \tfrac12 D\big(z(\theta_i)\big)\cos\theta_i, \\
y(\theta_i) &= \tfrac12 D\big(z(\theta_i)\big)\sin\theta_i, \\
z(\theta_i) &= z_i .
\end{aligned}
\label{eq:helix}
\end{equation}
For a constant-diameter cylindrical spring, $D(z)\equiv D$ and
Eq.~\eqref{eq:helix} reduces to the standard circular helix; for a
variable-geometry spring, $D(z)$ modulates the local helix radius,
producing conical, barrel, or otherwise shaped centerlines while the
pitch/angular relation $z(\theta)$ is inherited from the underlying
wire-length development.

\section{Summary of the computational pipeline}

The full property-calculation pipeline composes the results above in the
following order:
\begin{enumerate}
\item Compute base linear-spring properties (stiffness $k$, spring
index, wire length) from the underlying variable-lineal model.
\item Evaluate the Wahl factor and category at the reference mid-travel
position $z = L_0/2$ (Eq.~\eqref{eq:wahl}).
\item Determine the solid length $L_s$ via the telescoping test
(Eqs.~\eqref{eq:telescope-criterion}--\eqref{eq:solid-length}).
\item For each requested position $z_j \ge L_s$: obtain $F(z_j)$ from
the progressive-compression curve (Eqs.~\eqref{eq:incremental-force}--
\eqref{eq:load-at-position}), then $\tau(z_j)$
(Eq.~\eqref{eq:stress}) and $D_{\mathrm{out}}(z_j),D_{\mathrm{in}}(z_j)$
(Eqs.~\eqref{eq:outer-diameter}--\eqref{eq:inner-diameter}).
\item Extract $\tau_{\max},\tau_{\min}$ over all recorded positions and
evaluate the Goodman safety criterion
(Eqs.~\eqref{eq:mean-alt}--\eqref{eq:safety-factor}).
\end{enumerate}

\section{Conclusion}

The model presented here reformulates a variable-geometry compression
spring as: a position-indexed geometric field $D(z)$; a pointwise
Wahl-corrected torsion law (Eqs.~\eqref{eq:wahl},~\eqref{eq:stress}); a
combinatorial/geometric criterion selecting between two closed-form
solid-length regimes (Eq.~\eqref{eq:solid-length}); a discretized,
contact-aware load--deflection map replacing the classical linear spring
law (Eqs.~\eqref{eq:incremental-force}--\eqref{eq:interp}); and a
standard Goodman fatigue criterion applied to the resulting stress
extremes (Eq.~\eqref{eq:goodman}). Together these equations reproduce,
in closed mathematical form, the quantities computed and plotted by the
reference implementation.

\end{document}
